\chapter{内核内存分配器}

\section{内核分配面对的对象与上下文}

虚拟内存章节讲页表和地址空间，本章看内核自己怎样管理内存。内核需要分配页表页、进程结构、文件对象、socket buffer、目录项缓存、驱动 DMA buffer、日志记录和临时缓冲区。这些对象大小不同、生命周期不同、对齐要求不同，有些可以睡眠等待内存，有些在中断上下文不能睡眠。

先设想内核只提供“分配若干完整物理页”这一种操作。页表页和大块连续缓冲区能够直接使用，
但一个只有数百字节的文件对象也会独占整页；成千上万个小对象会把大部分内存留在页内空白中。
如果反过来只提供任意字节大小的通用堆，虽然小对象浪费较少，却很难在任意时刻取得满足页
对齐、物理连续和特定内存区域要求的大块空间。长期的任意切分还会使空闲区逐渐零散。

因此，内核必须把两个不同问题分开。底层先负责“哪些完整物理页可用，以及怎样取得连续的
$2^k$ 个页”；上层再负责“怎样把已经取得的页切成大量同类型对象，并快速复用这些对象槽”。
前一层通常由 buddy allocator 承担，后一层通常由 slab 或 SLUB 承担；通用小对象接口再把
不同请求大小映射到若干对象类别。专用的 task、inode 和 dentry 缓存位于更上层。

这种分层不是术语上的分类，而是因为单一分配器无法同时高效处理连续页、大量小对象、
对象初始化、缓存局部性和调试信息。后文先推导页级拆分与合并，再说明对象缓存为何建立在
页分配器之上。

\section{从物理页到内核对象的分配过程}

\subsection{页框为什么还要分区管理}

即使所有物理页大小相同，它们也不一定能够互换。有些设备只能访问特定地址范围；某些页
已经被内核固定，不能迁移；多处理器系统中，访问本地内存与远端内存的代价也可能不同。
如果分配器只维护一条全局空闲页链表，低地址页可能被普通请求提前耗尽，真正受地址限制的
设备随后即使看到大量空闲页也无法工作；所有 CPU 还会争用同一把链表锁。

因此，操作系统先把可分配页按可达范围和拓扑归入不同区域，再在每个区域内维护空闲状态。
分配请求不仅给出页数，还可能给出允许使用的区域、能否等待、能否回收以及是否允许移动
现有页面。所谓 zone、NUMA node 和迁移类型，分别是在表达这些限制；它们不是对物理内存
重复编号，而是分配决策所需的策略边界。

\begin{longtable}{|F{0.22\linewidth}|F{0.35\linewidth}|F{0.33\linewidth}|}
\toprule
\rowcolor{BookBlueTint}
约束 & 若忽略会发生什么 & 分配器需要保存的状态 \\
\midrule
地址可达范围 & 受限设备取得无法访问的页 & 页所属区域和请求允许区域 \\\tablerowrule
物理连续性 & 大页或设备缓冲无法建立 & 各阶连续空闲块 \\\tablerowrule
可移动性 & 内存压缩无法腾出连续区间 & movable、reclaimable、unmovable 分类 \\\tablerowrule
处理器局部性 & 频繁访问落到远端节点 & 节点归属、回退顺序和负载 \\\tablerowrule
执行上下文 & 中断路径进入睡眠或递归回收 & 是否允许等待、I/O 与文件系统回收 \\
\bottomrule
\end{longtable}

\subsection{连续页为何使用二次幂分组}

设一个区域只有 64 页，当前全部空闲。若把每一种可能长度的连续区间都登记起来，区间彼此
重叠，一次分配会使大量记录失效。若只逐页登记，请求 8 个连续页时又要扫描并验证相邻关系。
一种更容易维护的折中是只登记长度为 $1,2,4,8,\ldots$ 页且满足相应对齐的空闲块。

现在请求 8 页。分配器从 64 页块开始，先拆成两个 32 页块，保留其中一个；再把另一个拆成
两个 16 页块；最后拆成两个 8 页块并返回一个。每次拆分产生的另一半与返回方向中的块具有
确定关系，释放时可以通过块号计算出它的配对块。这种成对拆分与合并关系就是 buddy 名称的
来源。

\begin{lstlisting}
64 pages free
  -> 32 used for search + 32 free
  -> 16 used for search + 16 free + 32 free
  -> 8 allocated + 8 free + 16 free + 32 free

free the allocated 8 pages:
  8 + its free 8-page buddy -> 16
  16 + its free 16-page buddy -> 32
  32 + its free 32-page buddy -> 64
\end{lstlisting}

合并成立有两个条件：配对块当前必须空闲，而且必须处在同一阶。一个相邻的 8 页区间如果
属于另一个 16 页父块，就不是当前块的 buddy，不能仅凭地址相邻进行合并。这个规则使分配和
释放不需要扫描所有空闲区间，却也规定了请求大小要向上取整到二次幂。

若请求 9 页，分配器必须返回 16 页，其中 7 页暂时不能给其他请求使用，这称为内部碎片。
若系统总计有 32 个空闲页，但它们分散成许多不相邻单页，请求 8 个连续页仍然会失败，这称为
外部碎片。总空闲量只能说明容量，不能说明最大连续块；这正是高阶分配失败时不能只看剩余
内存百分比的原因。

\subsection{小对象为何从已分配页中复用}

接着考虑大小为 256 字节的 inode 内存对象。在 4096 字节页中，扣除少量管理信息后可以排列
多个同样大小的槽位。第一次请求需要向页分配器取得一页并建立槽位表，后续请求只需从空闲
槽链表取出一个；释放对象也只归还槽位。当整页所有槽位都空闲时，该页才有机会退回页分配器。

这层对象缓存不仅节约空间，还可以保留类型信息。分配器知道该缓存中的槽位都将承载同一种
对象，因而能够统一处理对齐、构造、统计、红区与释放后填充值。缓存通常把 slab 分成三类：
仍有空槽的 partial slab、已经用满的 full slab，以及完全空闲的 empty slab。分配优先使用
partial slab，避免在许多页上各留下少量无法归还的对象。

对象释放后不能立刻假定其所有引用都消失。引用计数、RCU 宽限期、设备完成和异步工作都可能
延长真实生命周期。分配器只负责槽位何时可以复用；对象所属子系统必须先证明再没有合法访问者。
将“从全局索引移除”和“把内存交回分配器”混成一步，是释放后使用错误的重要来源。

\section{页框区域、连续性与碎片}

实践先写对象分类表：

\begin{longtable}{|p{0.22\linewidth}|p{0.28\linewidth}|p{0.36\linewidth}|}
\toprule
对象 & 分配来源 & 约束 \\
\midrule
页表页 & buddy order-0 页 & 页对齐，通常不可移动 \\\tablerowrule
DMA buffer & 连续物理页或 IOMMU 映射 & 设备可访问，对齐和边界限制 \\\tablerowrule
task 结构 & slab cache & 高频分配释放，需要初始化 \\\tablerowrule
socket buffer & slab + page fragment & 网络吞吐和生命周期复杂 \\\tablerowrule
日志缓冲 & kmalloc 或专用池 & 失败路径也要可用 \\
\bottomrule
\end{longtable}

然后为每类对象写：是否可睡眠、是否可回收、是否可移动、是否需要清零、是否可能从中断上下文分配。

\section{buddy、slab 与分配上下文}

\subsection{buddy allocator}

buddy allocator 把空闲内存按 order 管理，order=k 表示 \code{2^k} 个连续页。分配时找最小足够 order；没有就拆更大的块。释放时尝试和 buddy 合并。

\begin{lstlisting}
alloc(order):
  for k in order..MAX_ORDER:
    if free[k] not empty:
      block = free[k].pop()
      while k > order:
        k -= 1
        buddy = split(block, k)
        free[k].push(buddy)
      return block
  return null

free(block, order):
  while order < MAX_ORDER:
    buddy = find_buddy(block, order)
    if buddy not free: break
    remove buddy from free[order]
    block = merge(block, buddy)
    order += 1
  free[order].push(block)
\end{lstlisting}

buddy 的优势是能管理连续页，释放时能合并。问题是内部碎片：请求 9 页必须分配 16 页。长期运行后，高 order 连续块也可能稀缺。

\subsection{slab cache}

slab 分配器为固定大小对象建立缓存。一个 slab 通常由一页或多页组成，切成多个对象槽。分配就是从 freelist 取一个槽；释放就是放回 freelist。对象构造可以在 slab 创建时完成，减少重复初始化成本。

\begin{lstlisting}
slab_alloc(cache):
  if cache.partial not empty:
    obj = pop_free_object(cache.partial.first)
    return obj
  slab = allocate_new_slab(cache)
  cache.partial.push(slab)
  return pop_free_object(slab)
\end{lstlisting}

slab 的分配释放接近 \code{O(1)}。它适合大量同类型对象，例如 inode、dentry、task。实现要处理 per-CPU cache，以减少多核锁竞争；还要处理调试模式下的红区、poison、double free 检测。

\subsection{kmalloc 和大小类别}

通用 \code{kmalloc} 通常把请求大小映射到一组 size class，例如 32、64、128、256 字节。每个 size class 背后是一个 slab cache。这样可以用固定对象缓存服务变长请求。

\begin{lstlisting}
kmalloc(size):
  class = size_to_class(size)
  return slab_alloc(kmalloc_cache[class])
\end{lstlisting}

size class 的代价是内部碎片。请求 129 字节可能分到 256 字节槽。分配器设计要在碎片、速度和缓存局部性之间折中。

\subsection{GFP 标志和上下文}

内核分配 API 通常带标志，说明是否允许睡眠、是否允许 I/O 回收、是否要求 DMA 区域等。中断上下文不能睡眠，所以只能使用不会阻塞的分配；普通进程上下文可以等待回收。

\begin{lstlisting}
if in_interrupt():
  page = alloc_page(GFP_ATOMIC)
else:
  page = alloc_page(GFP_KERNEL)
\end{lstlisting}

这解释了为什么内核代码不能随便分配内存。分配失败不是理论问题，尤其在中断、持锁、内存压力和回收路径中。高质量代码必须定义失败策略：返回错误、丢包、使用预留池、降级还是 panic。

\subsection{per-CPU 缓存}

多核系统中，所有小对象分配都争同一把 slab 锁会很慢。per-CPU cache 让每个 CPU 维护少量本地空闲对象，常见分配释放不需要全局锁。

\begin{lstlisting}
kmem_cache_alloc(cache):
  cpu_cache = cache.percpu[current_cpu()]
  if cpu_cache.freelist not empty:
    return cpu_cache.freelist.pop()
  refill_from_global(cache, cpu_cache)
  return cpu_cache.freelist.pop()

kmem_cache_free(cache, obj):
  cpu_cache = cache.percpu[current_cpu()]
  cpu_cache.freelist.push(obj)
  if cpu_cache.count > HIGH_WATERMARK:
    drain_to_global(cache, cpu_cache)
\end{lstlisting}

per-CPU 缓存降低锁竞争，但会增加内存滞留：每个 CPU 都可能囤积一批对象。高质量实现要有高低水位、CPU 下线 drain、内存压力回收和统计对账。

\subsection{内存调试与防御}

分配器是许多内核漏洞的放大器。越界写、use-after-free、double free 和未初始化读取都可能破坏任意内核对象。调试构建应加入红区、poison、隔离队列和调用栈记录。

\begin{lstlisting}
debug_free(obj):
  if obj.magic != LIVE_MAGIC:
    panic("double free or corrupted object")
  fill(obj.memory, POISON_FREE)
  obj.magic = FREE_MAGIC
  quarantine_push(obj)

debug_alloc(cache):
  obj = real_alloc(cache)
  fill(obj.memory, POISON_ALLOC)
  obj.magic = LIVE_MAGIC
  return obj
\end{lstlisting}

隔离队列会延迟复用刚释放对象，使 use-after-free 更容易在测试中暴露。代价是内存占用和性能下降，所以通常只在调试或安全强化配置中开启。

\section{泄漏、越界与对象生命周期}

\begin{itemize}
\tightlist
\item buddy 分配释放后，空闲页总数守恒。
\item 分裂和合并后，不存在重叠空闲块。
\item slab 分配对象满足对齐，释放后能复用。
\item double free、越界写和 use-after-free 在调试模式能被检测。
\item GFP\_ATOMIC 路径不能睡眠，失败时有明确降级策略。
\item 内存压力下，分配器统计能解释碎片和失败原因。
\end{itemize}

\section{SLUB、内存压缩与资源归属}

\subsection{buddy 的 free\_area 和迁移类型}

Linux buddy 不只是按 order 分链表，还按迁移类型分组，例如 movable、unmovable、reclaimable、CMA 等。这样做是为了降低长期碎片：可移动页尽量放在一起，未来需要高 order 连续页时更容易通过 compaction 得到。

\begin{lstlisting}
struct free_area {
  list_head free_list[MIGRATE_TYPES];
  unsigned long nr_free;
};
\end{lstlisting}

页分配路径必须同时考虑空闲链表、水位线、内存区域和迁移类型。一次 order-9 分配失败，
不一定表示系统没有 512 页空闲，也可能只是没有连续的 512 页。

\subsection{SLAB、SLUB、SLOB 的取舍}

SLAB 强调对象缓存和较复杂的队列；SLUB 简化元数据，成为常见默认；SLOB 面向极小系统，
结构简单但伸缩性较弱。理解 SLUB 时，应把对象缓存、每 CPU 空闲链表、部分使用的 slab，
以及调试用红区和填充值联系起来。

\begin{lstlisting}
kmem_cache_alloc(cache):
  if current_cpu.freelist:
    return pop(current_cpu.freelist)
  if cache.node.partial:
    slab = take_partial(cache.node)
    load_cpu_freelist(slab)
    return pop(current_cpu.freelist)
  slab = new_slab_from_buddy()
  return first_object(slab)
\end{lstlisting}

SLUB 的快路径很短，但 debug 选项会显著增加检查成本。教材中讲“分配是 O(1)”时，要同时说明这是平均快路径；慢路径可能进入 buddy、回收、compaction 或 OOM。

\subsection{kmalloc size class 与内部碎片}

\code{kmalloc(130)} 可能进入 192 或 256 字节 size class，取决于架构和配置。向上取整带来内部碎片，但换来固定大小缓存的速度和对齐。对象若频繁分配，最好建立专用 \code{kmem\_cache}，可以提供构造函数、调试名和更好统计。

\subsection{GFP\_NOIO、GFP\_NOFS 和回收递归}

\code{GFP\_KERNEL} 允许睡眠和回收；\code{GFP\_ATOMIC} 用于中断或持自旋锁路径；\code{GFP\_NOIO} 禁止通过 I/O 回收；\code{GFP\_NOFS} 禁止进入文件系统回收路径。后两者用于避免递归死锁。

\begin{lstlisting}
filesystem_write_path():
  lock(inode)
  buf = kmalloc(size, GFP_NOFS)
  // avoid reclaim entering filesystem and taking inode locks again
\end{lstlisting}

若文件系统持有 inode 锁时用 \code{GFP\_KERNEL} 分配，内存压力下回收可能进入同一文件系统并等待 inode 锁，形成递归死锁。GFP flag 是锁和回收的边界，不是性能装饰。

\subsection{vmalloc 与物理不连续}

\code{kmalloc} 通常返回物理连续内存，适合 DMA 或小对象；\code{vmalloc} 返回虚拟连续但物理页可不连续的内存，适合大缓冲或模块映射。代价是需要建立页表，TLB 局部性更差，不能直接给不支持 scatter-gather 的设备 DMA。

\begin{lstlisting}
vmalloc(size):
  area = allocate_virtual_area(size)
  pages = alloc_individual_pages(npages)
  map_pages(area, pages)
  return area.addr
\end{lstlisting}

\subsection{kmemleak、KASAN 和 KFENCE}

内存泄漏和 use-after-free 很难靠肉眼找。kmemleak 通过扫描可达指针寻找未释放对象；KASAN 使用 shadow memory 检测越界和释放后访问；KFENCE 用低频保护页检测堆错误，开销比 KASAN 低但覆盖率也低。

\begin{longtable}{|p{0.18\linewidth}|p{0.36\linewidth}|p{0.34\linewidth}|}
\toprule
工具 & 能发现 & 代价 \\
\midrule
kmemleak & 长期未释放对象 & 扫描和误报，需要调试配置 \\\tablerowrule
KASAN & 越界、use-after-free & 内存和性能开销高 \\\tablerowrule
KFENCE & 采样式堆越界和 UAF & 开销低，非确定覆盖 \\
\bottomrule
\end{longtable}

\subsection{内存碎片、compaction 和 CMA}

内存碎片分为内部碎片和外部碎片。内部碎片来自 size class 向上取整；外部碎片来自空闲页分散，无法满足高 order 连续页。compaction 通过迁移可移动页，把空闲页聚合成连续块。CMA 预留可迁移区域，给需要连续物理内存的设备使用。

\begin{lstlisting}
compact_zone(zone):
  free_scanner = end_of_zone
  migrate_scanner = start_of_zone
  while scanners_not_met:
    page = find_movable_page(migrate_scanner)
    free = find_free_page(free_scanner)
    migrate(page, free)
\end{lstlisting}

并非所有页都能迁移。页表页、驱动 pin 住的 DMA 页、mlock 页和某些内核对象不可移动。高 order 分配失败时，诊断要看总空闲页、最大连续块、不可移动页比例和 compaction 失败原因。

\subsection{内存控制组和分配归属}

cgroup v2 可以限制一组进程的内存。分配器需要把许多用户可归因的内存记到 memcg，例如 page cache、匿名页、部分 slab 对象。这样容器内泄漏不会无限吃掉宿主内存。

\begin{lstlisting}
charge_memcg(page, current.memcg):
  if memcg.current + page.size > memcg.max:
    reclaim_memcg(memcg)
    if still over limit:
      trigger_memcg_oom(memcg)
  memcg.current += page.size
\end{lstlisting}

memcg OOM 和全局 OOM 必须区分。前者应限制在相应资源组内，后者说明整机资源不足。
判断时要同时核对资源组当前用量、内存事件和最终被选择终止的任务。

\subsection{对象生命周期对账}

内核内存质量不只看分配成功，还要能证明对象生命周期闭合。每类对象应有 alloc/free 计数、当前活跃数、高水位和失败次数。对复杂对象，还应记录 owner 和引用来源。

\begin{lstlisting}
object_accounting:
  alloc_count++
  active_count++
  high_watermark = max(high_watermark, active_count)

free:
  active_count--
  free_count++
\end{lstlisting}

测试结束时 active count 不为零，就是泄漏线索；free count 大于 alloc count，说明 double free 或计数损坏；高水位异常说明负载下对象滞留。生产系统不一定记录所有栈，但调试构建应支持按对象类型打开详细追踪。
